% 1A mock exam (AP-style). Build: python3 units/ecology/study-guides/build.py 1A-exam
% Macros: latex/exam.sty. Figures: \drawfig/\examfig (figures/) and \bookcrop (cropped from a book by build.py).
\documentclass[11pt,letterpaper]{article}
\newcommand{\sgroot}{../}
% Shared preamble for the study guides and mock exams (built by ../build.py).
% \sgroot is the path from the document's folder to study-guides/ (empty for guides, ../ for exams).
\providecommand{\sgroot}{}
\usepackage[margin=0.8in,headheight=15pt]{geometry}
\usepackage{fontspec}
\setmainfont{texgyrepagella}[Extension=.otf,UprightFont=*-regular,BoldFont=*-bold,ItalicFont=*-italic,BoldItalicFont=*-bolditalic]
\setsansfont{texgyreheros}[Extension=.otf,UprightFont=*-regular,BoldFont=*-bold,ItalicFont=*-italic,BoldItalicFont=*-bolditalic]
\usepackage{amsmath,amssymb,booktabs,longtable,array,calc,enumitem,xcolor,graphicx}
\usepackage{tikz}
\usepackage{pgfplots}\pgfplotsset{compat=1.18}
\usetikzlibrary{calc,arrows.meta,positioning}
\usepackage[most]{tcolorbox}
\usepackage{fancyhdr,titlesec,needspace,etoolbox,environ}
\usepackage{hyperref}

\definecolor{forest}{HTML}{24594B}
\definecolor{pale}{HTML}{EFF5F1}
\definecolor{keyc}{HTML}{C98A0B}   % key idea: amber
\definecolor{vocabc}{HTML}{2F6FA8} % vocabulary: blue
\definecolor{tipc}{HTML}{3F8A3A}   % make it make sense: green
\definecolor{watchc}{HTML}{B8452F} % watch out: red
\definecolor{linkc}{HTML}{7B4C9A}  % connection: purple
\hypersetup{colorlinks=true,linkcolor=forest,urlcolor=vocabc}

\pagestyle{fancy}\fancyhf{}
\renewcommand{\headrulewidth}{0pt}
\fancyfoot[C]{\small\sffamily\thepage}
\setlength{\parindent}{0pt}\setlength{\parskip}{5pt}
\setlength{\emergencystretch}{2em}
\setlist[itemize]{leftmargin=1.3em,itemsep=2pt,topsep=2pt}
\setlist[enumerate]{leftmargin=1.6em,itemsep=2pt,topsep=2pt}
\renewcommand{\arraystretch}{1.2}
\providecommand{\tightlist}{\setlength{\itemsep}{0pt}\setlength{\parskip}{0pt}}
\providecommand{\pandocbounded}[1]{#1}
\def\LTcaptype{none}

\titleformat{\section}{\Large\sffamily\bfseries\color{forest}}{}{0pt}{}[{\color{forest}\titlerule[0.8pt]}]
\titleformat{\subsection}{\large\sffamily\bfseries\color{forest}}{}{0pt}{}
\titleformat{\subsubsection}{\normalsize\sffamily\bfseries}{}{0pt}{}
\titlespacing*{\section}{0pt}{14pt}{8pt}
\titlespacing*{\subsection}{0pt}{12pt}{4pt}
\setcounter{secnumdepth}{0}

% Block quotes (textbook excerpts) as shaded boxes.
\renewenvironment{quote}{\begin{tcolorbox}[breakable,enhanced,colback=pale,colframe=forest,boxrule=0pt,leftrule=2.5pt,arc=0pt,
  left=8pt,right=8pt,top=4pt,bottom=4pt,before skip=4pt,after skip=6pt]\small}{\end{tcolorbox}}

% ------------------------------------------------------------------ annotated textbook pages
% Landscape letter page: the scanned page on the left (full height), notes on the right.
\newlength{\portraitW}\setlength{\portraitW}{8.5in}
\newlength{\portraitH}\setlength{\portraitH}{11in}
\def\pgmargin{21.6}  % pt, top/left margin around the scan
\def\pgheight{568.8} % pt, printed height of the scan (8.5in - 2 margins)
\newcommand{\setpagesize}[2]{%
  \global\pdfpagewidth=#1\global\pdfpageheight=#2%
  \global\paperwidth=#1\global\paperheight=#2}

\newcommand{\circlednum}[2]{\tikz[baseline=(c.base)]\node[circle,fill=#1,text=white,font=\scriptsize\sffamily\bfseries,inner sep=1.2pt,minimum size=11pt](c){#2};}

% \textbookpage{image}{page width pt}{page height pt}{lesson}{book page}{intro}{highlights}{notes}
% Notes (\addnote) are typeset into boxes first, then laid out: each sits level with its highlight
% when there is room, is pushed down to avoid the note above, and is pulled back up from the page bottom.
\newcount\nnotes
\newbox\hdrbox
\newdimen\prevbottom \newdimen\limitdim \newdimen\tmpdim
\def\notegap{7pt}
\newcommand{\addnote}[5]{% color, number, label, anchor (scan pt, <0 = none), text
  \global\advance\nnotes by 1
  \expandafter\ifx\csname nb@\the\nnotes\endcsname\relax\expandafter\newbox\csname nb@\the\nnotes\endcsname\fi
  \global\expandafter\setbox\csname nb@\the\nnotes\endcsname=\hbox{\tikz\node[notebox=#1]{\notehead{#1}{#2}{#3}#5};}%
  \expandafter\xdef\csname na@\the\nnotes\endcsname{#4}%
  \expandafter\xdef\csname nc@\the\nnotes\endcsname{#1}}
\newcommand{\nb}[1]{\csname nb@#1\endcsname}
\newcommand{\layoutnotes}{%
  \global\prevbottom=\dimexpr\pgmargin pt+\ht\hdrbox+\dp\hdrbox\relax
  \ifnum\nnotes>0
    \foreach \k in {1,...,\the\nnotes}{%
      \edef\a{\csname na@\k\endcsname}%
      \tmpdim=\dimexpr\prevbottom+\notegap\relax
      \ifdim \a pt<0pt\else
        \pgfmathsetlength{\dimen0}{\pgmargin+\a*\s-7}%
        \ifdim\dimen0>\tmpdim \tmpdim=\dimen0 \fi
      \fi
      \expandafter\xdef\csname ny@\k\endcsname{\the\tmpdim}%
      \global\prevbottom=\dimexpr\tmpdim+\ht\nb{\k}+\dp\nb{\k}\relax}%
    \global\limitdim=\dimexpr612pt-32pt\relax
    \foreach \k in {\the\nnotes,...,1}{%
      \tmpdim=\csname ny@\k\endcsname
      \ifdim\dimexpr\tmpdim+\ht\nb{\k}+\dp\nb{\k}\relax>\limitdim
        \tmpdim=\dimexpr\limitdim-\ht\nb{\k}-\dp\nb{\k}\relax
        \expandafter\xdef\csname ny@\k\endcsname{\the\tmpdim}%
      \fi
      \global\limitdim=\dimexpr\tmpdim-\notegap\relax}%
    \ifdim\limitdim<\dimexpr\pgmargin pt+\ht\hdrbox+\dp\hdrbox-\notegap\relax
      \typeout{WARNING: notes overflow on textbook page \currentbookpage}\fi
  \fi}
\newcommand{\drawnotes}{%
  \ifnum\nnotes>0
    \foreach \k in {1,...,\the\nnotes}{%
      \edef\a{\csname na@\k\endcsname}\edef\c{\csname nc@\k\endcsname}\edef\y{\csname ny@\k\endcsname}%
      \node[anchor=north west,inner sep=0pt] at ($(O)+(\notex pt,-\y)$) {\usebox{\nb{\k}}};
      \ifdim \a pt<0pt\else
        \draw[\c,line width=0.6pt] ($(scan.north east)+(0,{-\a*\s pt})$) -- ++(6pt,0) -- ($(O)+(\notex pt-1pt,-\y-7pt)$);
      \fi}%
  \fi}
\newcommand{\textbookpage}[8]{%
  \clearpage\setpagesize{11in}{8.5in}\thispagestyle{empty}%
  \def\currentbookpage{#5}%
  \pgfmathsetmacro{\s}{\pgheight/#3}%
  \pgfmathsetmacro{\pgright}{\pgmargin+#2*\s}%
  \pgfmathsetmacro{\notex}{\pgright+20}%
  \pgfmathsetmacro{\notew}{792-\notex-\pgmargin}%
  \pgfmathsetmacro{\notetw}{\notew-10}%
  \setbox\hdrbox=\vbox{\hsize=\notew pt\parskip=0pt
    {\raggedright\sffamily\bfseries\color{forest}#4\par}\vspace{1pt}
    {\sffamily\footnotesize\color{black!55}Miller \& Levine Biology, textbook page #5\par}\vspace{-1pt}
    {\color{forest}\rule{\notew pt}{0.6pt}}\par\vspace{2pt}
    \ifx\relax#6\relax\else{\small\raggedright #6\par}\fi}%
  \global\nnotes=0 #8\layoutnotes
  \null
  \begin{tikzpicture}[remember picture,overlay]
    \coordinate (O) at (current page.north west);
    \node[anchor=north west,inner sep=0pt,draw=black!25] (scan) at ($(O)+(\pgmargin pt,-\pgmargin pt)$)
      {\includegraphics[height=\pgheight pt]{#1}};
    \node[anchor=north west,inner sep=0pt] at ($(O)+(\notex pt,-\pgmargin pt)$) {\usebox{\hdrbox}};
    \node[anchor=south east,font=\scriptsize\sffamily,text=black!50,inner sep=0pt] at ($(O)+(792pt-\pgmargin pt,-603pt)$)
      {Part 1 · the textbook, annotated \quad·\quad Miller \& Levine Biology (2019), p.\,#5 \quad·\quad page \thepage};
    #7
    \drawnotes
  \end{tikzpicture}%
  \clearpage\setpagesize{\portraitW}{\portraitH}}

% Highlight one text line on the scan: \hl{color}{x0}{y0}{x1}{y1} in scan points (y down).
\newcommand{\hl}[5]{%
  \fill[#1,opacity=0.22] ($(scan.north west)+({(#2-1.5)*\s pt},{-(#3-1)*\s pt})$) rectangle
                         ($(scan.north west)+({(#4+1.5)*\s pt},{-(#5+1)*\s pt})$);}
% Numbered badge just left of a highlight.
\newcommand{\badge}[4]{%
  \node[circle,fill=#1,text=white,font=\scriptsize\sffamily\bfseries,inner sep=1pt,minimum size=10pt]
    at ($(scan.north west)+({#3*\s pt-7pt},{-#4*\s pt})$) {#2};}
\newcommand{\notehead}[3]{{\sffamily\bfseries\footnotesize\circlednum{#1}{#2}\ \color{#1}\MakeUppercase{#3}}\\[1pt]}
\tikzset{notebox/.style={text width=\notetw pt,inner sep=4pt,fill=#1!7,draw=#1!55,line width=0.4pt,
  font=\small,align=left,rounded corners=1.5pt}}

% ------------------------------------------------------------------ figures in the guide text
% Caption used under both book figures and new diagrams.
\newcommand{\figcaption}[2]{\par\smallskip{\small\sffamily\raggedright #1\par}%
  \ifx\relax#2\relax\else{\scriptsize\sffamily\color{black!55}Source: #2\par}\fi}
% \bookfig{png}{width fraction}{caption}{source}: a figure cropped from one of the books.
\newcommand{\bookfig}[4]{\par\Needspace{8\baselineskip}\medskip
  \begin{minipage}{\linewidth}\centering
    \fcolorbox{black!15}{white}{\includegraphics[width=#2\linewidth]{#1}}
    \figcaption{#3}{#4}
  \end{minipage}\par\medskip}
% Several book figures side by side: \bookfigrow{\bookfigcol{w}{png}{height}{caption}{source}\hfill ...}
\newcommand{\bookfigcol}[5]{\begin{minipage}[t]{#1\linewidth}\centering
    \fcolorbox{black!15}{white}{\includegraphics[height=#3,width=0.95\linewidth,keepaspectratio]{#2}}
    \figcaption{#4}{#5}\end{minipage}}
\newcommand{\bookfigrow}[1]{\par\Needspace{10\baselineskip}\medskip\noindent#1\par\medskip}
% New diagrams drawn for the guide: \begin{guidefig}...\end{guidefig}{caption}
\newenvironment{guidefig}[1]{\par\Needspace{10\baselineskip}\medskip\begin{minipage}{\linewidth}\centering\def\guidefigcap{#1}}
  {\figcaption{\guidefigcap}{}\end{minipage}\par\medskip}
\tikzset{
  gbox/.style={draw=#1,fill=#1!8,rounded corners=3pt,align=center,font=\small\sffamily,inner sep=5pt},
  garr/.style={-{Stealth[length=2.4mm]},line width=1pt,draw=#1},
  glab/.style={font=\scriptsize\sffamily,fill=white,inner sep=1.5pt,align=center},
}

% ------------------------------------------------------------------ figures shared by guides and exams
% \drawfig{name}: a TikZ/pgfplots figure in figures/name.tex
\newcommand{\drawfig}[1]{\input{\sgroot figures/#1}}
% \bookcrop[width]{id}{book}{page}{x0 y0 x1 y1}{label}{caption}: a figure cropped from a book page.
% book: C = Campbell Biology, ML = Miller & Levine Biology; page = printed page number;
% crop box in points from the page's top-left corner. build.py makes build/figs/<id>.png from these arguments.
\newcommand{\bookname}[1]{\ifstrequal{#1}{C}{Campbell Biology}{Miller \& Levine Biology}}
\newcommand{\bookcrop}[7][0.7]{\bookfig{\sgroot build/figs/#2.png}{#1}{#7}{\bookname{#3}, #6, p.\,#4}}

\usepackage{../latex/exam}
\hypersetup{pdftitle={1A Mock Exam},pdfauthor={Prepared for Shawn}}

\begin{document}
\examtitle{1A}{1A Mock Exam: Science Practices and Scientific Reasoning}
\textbf{Format.} This practice exam follows the AP Biology exam format: data-based multiple choice, then free response. It covers learning target 1A: observation and inference, types of data, inductive and deductive reasoning, experimental design, and how the scientific community checks ideas.

\begin{center}
\begin{tabular}{@{}lll@{}}
\toprule
Section & Questions & Suggested time \\
\midrule
I. Multiple choice & 16 & 25 minutes \\
II. Free response & 4 (one long, three short) & 45 minutes \\
\bottomrule
\end{tabular}
\end{center}

Answer in full sentences in Section II. Graphs and tables marked \textbf{hypothetical} were made up for practice. The answer key, with explanations and scoring guides, starts after Section II. Don't look until you've finished.

\clearpage

\examsection{Section I. Multiple choice}

Choose the \textbf{one} best answer for each question.

\questiongroup{Questions 1–3 refer to the field notes below.}

A student visits a pond and writes these notes:

\begin{fieldnotes}
  \item The water is green and cloudy.
  \item There are 14 frogs on the north bank and 3 frogs on the south bank.
  \item The frogs prefer the north bank because it gets more sun.
  \item The water is 18\,°C at the north bank and 15\,°C at the south bank.
  \item The pond is polluted.
\end{fieldnotes}

\begin{mcq}{1}
Which of the student's statements are inferences rather than observations?
\begin{choices}
  \item 1 and 2
  \item 2 and 4
  \item 3 and 5
  \item 1 and 4
\end{choices}
\end{mcq}

\begin{mcq}{2}
Which statements record quantitative data?
\begin{choices}
  \item 1 and 5
  \item 2 and 4
  \item 1 and 3
  \item 3 and 4
\end{choices}
\end{mcq}

\begin{mcq}{3}
Which of the following is a testable scientific hypothesis based on the notes?
\begin{choices}
  \item Frogs are happier on the north bank than on the south bank.
  \item If frogs choose warmer water, then more frogs will gather in a section of the pond that is warmed than in a similar section that is not warmed.
  \item The pond is green, so it must be polluted.
  \item The frogs on the north bank should be protected from people.
\end{choices}
\end{mcq}

\questiongroup{Questions 4–7 refer to the following experiment.}

Researchers wanted to know whether fur color that matches the surroundings protects mice from predators. They placed light-colored and dark-colored model mice in two habitats: a beach with light sand and an inland area with dark soil. The next morning they recorded the percentage of models in each group that had been attacked by predators.

\bookcrop[0.68]{c-mice-results}{C}{20}{58 290 380 482}{Figure 1.25}
  {Percentage of light and dark mouse models attacked by predators in a beach habitat and an inland habitat.}

\begin{mcq}{4}
Which is the independent variable in this experiment?
\begin{choices}
  \item The percentage of models attacked
  \item The color of the models
  \item The kind of predator that attacked
  \item The time of the next morning
\end{choices}
\end{mcq}

\begin{mcq}{5}
In the beach habitat, dark models were attacked about how many times as often as light models?
\begin{choices}
  \item 1.5 times
  \item 2.6 times
  \item 4.5 times
  \item 45 times
\end{choices}
\end{mcq}

\begin{mcq}{6}
Which claim is best supported by the data?
\begin{choices}
  \item Predators attack dark mice more often than light mice in every habitat.
  \item Models that did not match their habitat were attacked more often than models that matched it.
  \item The results prove that light fur evolved in beach mice because of predators.
  \item Dark models were attacked more often than light models overall.
\end{choices}
\end{mcq}

\begin{mcq}{7}
Why did the researchers test the models in both habitats instead of only on the beach?
\begin{choices}
  \item To increase the number of predators in the study
  \item To show that the effect depends on matching the background, not on one color always being safer
  \item To remove the independent variable from the experiment
  \item To prove that the hypothesis is true
\end{choices}
\end{mcq}

\questiongroup{Questions 8–9 refer to types of reasoning.}

\begin{mcq}{8}
A class tested 40 samples of water from different ponds. Every sample contained bacteria, so the class concluded, ``All pond water contains bacteria.'' Which statement correctly describes this reasoning?
\begin{choices}
  \item It is deductive, so the conclusion is guaranteed to be true.
  \item It is inductive, so the conclusion is probable but could be overturned by new evidence.
  \item It is deductive, so the conclusion is only probably true.
  \item It is inductive, so the conclusion is guaranteed to be true.
\end{choices}
\end{mcq}

\begin{mcq}{9}
Which of the following is an example of deductive reasoning?
\begin{choices}
  \item All insects have six legs. A beetle is an insect. Therefore, a beetle has six legs.
  \item The sun has risen every morning so far, so it will rise tomorrow.
  \item The numbers 3, 6, 9, 12 are followed by 15, so the rule is ``add 3.''
  \item Every swan I have seen is white, so all swans are white.
\end{choices}
\end{mcq}

\questiongroup{Questions 10–11 refer to the graph below.}

A class played the Geotastic game. In each trial, students could ``advance'' (move) a different number of times before guessing their location. They predicted: \emph{if more advances are allowed, then guesses will be closer to the actual location.}

\examfig{exam-1A-geotastic}

\begin{mcq}{10}
Which statement best describes the results?
\begin{choices}
  \item As the number of advances increased, the mean distance increased, which refutes the prediction.
  \item As the number of advances increased, the mean distance decreased, which supports the prediction.
  \item The results prove that more observations always lead to a correct inference.
  \item The number of advances had no effect on the mean distance.
\end{choices}
\end{mcq}

\begin{mcq}{11}
Which change would make the class's conclusion \textbf{more reliable}?
\begin{choices}
  \item Using only one student's results
  \item Combining data from all class periods to increase the sample size
  \item Allowing students to search the internet for location clues
  \item Changing both the number of advances and the time limit in each trial
\end{choices}
\end{mcq}

\questiongroup{Questions 12–13 refer to the graph below.}

Researchers added different amounts of nitrogen fertilizer to plots of marsh grass and measured the grass after 8 weeks.

\examfig{exam-1A-nitrogen}

\begin{mcq}{12}
Which conclusion is best supported by the data?
\begin{choices}
  \item Grass height increased from 0 to 10 g/m\textsuperscript{2} of nitrogen, but the data do not show a clear difference between 10 and 20 g/m\textsuperscript{2}.
  \item Every increase in nitrogen caused a significant increase in grass height.
  \item Nitrogen had no effect on grass height.
  \item Grass given 20 g/m\textsuperscript{2} of nitrogen grew almost twice as tall as grass given 10 g/m\textsuperscript{2}.
\end{choices}
\end{mcq}

\begin{mcq}{13}
What is the main purpose of the plots that received 0 g/m\textsuperscript{2} of nitrogen?
\begin{choices}
  \item They are the dependent variable.
  \item They are a control group that shows how the grass grows without added nitrogen, for comparison.
  \item They make sure the hypothesis will be supported.
  \item They keep the amount of sunlight the same in all plots.
\end{choices}
\end{mcq}

\questiongroup{Questions 14–16}

\begin{mcq}{14}
Which statement about scientific theories is correct?
\begin{choices}
  \item A theory is an educated guess that has not yet been tested.
  \item A theory becomes a law once it has been proven.
  \item A theory is a well-tested explanation that ties together many observations, and it can still be revised if new evidence appears.
  \item Once a theory is accepted, scientists stop testing it.
\end{choices}
\end{mcq}

\begin{mcq}{15}
A research team publishes its results in a peer-reviewed journal. Which statement best describes what peer review does?
\begin{choices}
  \item It guarantees that the conclusions are correct.
  \item Other experts check the methods, data, and reasoning for mistakes, bias, or fraud before publication.
  \item It lets the researchers' friends approve the paper.
  \item It replaces the need for other scientists to replicate the study.
\end{choices}
\end{mcq}

\begin{mcq}{16}
A student scores 60\% on a practice test that does not count toward the grade. Based on the idea of deliberate practice, what is the most effective next step?
\begin{choices}
  \item Ignore the result, because practice tests don't affect the grade
  \item Reread all of the notes for the unit from the beginning
  \item Find the specific questions and skills that were missed, practice those skills, then check again with new questions
  \item Wait for the unit exam to find out what needs work
\end{choices}
\end{mcq}

\clearpage

\examsection{Section II. Free response}

\frq{Question 1 (long response, 10 points)}

A student wants to know whether the amount of light affects the growth of radish seedlings. She plants 40 seeds of the same variety in identical pots and divides them into four groups of 10. Each group gets a different number of hours of light per day. After 10 days she measures each seedling.

\begin{center}
\begin{tabular}{@{}lcccc@{}}
\toprule
Hours of light per day & 4 & 8 & 12 & 16 \\
\midrule
Mean height (cm) & 2.1 & 3.8 & 5.2 & 5.4 \\
$\pm$ 2 SE (cm) & 0.4 & 0.5 & 0.6 & 0.7 \\
\bottomrule
\end{tabular}
\end{center}

\datanote{Hypothetical data. n = 10 seedlings per group.}

\begin{parts}
  \item \textbf{Identify} the independent variable and the dependent variable in this experiment.
  \item \textbf{Identify} two variables that should be kept the same for all groups, and \textbf{explain} why keeping them the same matters.
  \item The student did not include a group grown in complete darkness. \textbf{Explain} whether the 4-hour group can serve as a comparison group, and \textbf{describe} one limitation of not having a 0-hour group.
  \item \textbf{Construct} an appropriate graph of the mean heights, with error bars, on the grid below. Label the axes.

  \graphgrid
  \item \textbf{Describe} the trend in the data. Using the error bars, \textbf{determine} whether the difference between the 12-hour and 16-hour groups is likely to be real.
  \item \textbf{Make a claim} about the effect of light on radish growth, and \textbf{support} it with evidence and reasoning (CER).
  \item \textbf{Predict} the mean height of seedlings grown with 20 hours of light per day, and \textbf{justify} your prediction.
\end{parts}

\frq{Question 2 (short response, 4 points)}

A biologist finds many dead fish in a river just downstream from a factory, but no dead fish upstream from the factory.

\begin{parts}
  \item \textbf{State} one observation and one inference from this description.
  \item The biologist concludes, ``The factory's wastewater killed the fish.'' \textbf{Identify} the type of reasoning (inductive or deductive) and \textbf{explain} why the conclusion is not certain.
  \item \textbf{Write} a testable hypothesis in ``if \dots{} then \dots{}'' form and \textbf{describe} the result that would support it.
  \item \textbf{Describe} one additional observation that would weaken the biologist's conclusion.
\end{parts}

\frq{Question 3 (short response, 4 points)}

Refer to the mouse model experiment in Questions 4–7 of Section I.

\begin{parts}
  \item \textbf{Identify} the control group in the inland habitat.
  \item \textbf{Calculate} the difference, in percentage points, between the attack rates of light and dark models in the inland habitat.
  \item The researchers used painted models instead of live mice. \textbf{Explain} one way this made the experiment better controlled.
  \item \textbf{Explain} how the results support the hypothesis that fur color that matches the surroundings protects mice from predators.
\end{parts}

\frq{Question 4 (short response, 4 points)}

A classmate says, ``This study was published in a science journal, so it's been proven true.''

\begin{parts}
  \item \textbf{Describe} what happens during peer review.
  \item \textbf{Explain} why being published does not mean a conclusion is proven.
  \item \textbf{Explain} how replication by other scientists makes a conclusion more trustworthy.
  \item \textbf{Explain} how formative feedback in your biology class is similar to peer review in science.
\end{parts}

\answerkey

\keyheading{Section I at a glance}

\begin{center}
\begin{tabular}{@{}lccccccc@{}}
\toprule
Q & Answer & Q & Answer & Q & Answer & Q & Answer \\
\midrule
1 & C & 5 & B & 9 & A & 13 & B \\
2 & B & 6 & B & 10 & B & 14 & C \\
3 & B & 7 & B & 11 & B & 15 & B \\
4 & B & 8 & B & 12 & A & 16 & C \\
\bottomrule
\end{tabular}
\end{center}

\textbf{Scoring guide.} 14–16 correct: ready. 11–13: review the explanations for the ones you missed. 10 or fewer: reread the study guide's Parts 2–3, then retake.

\keyheading{Section I explanations}

\answer{1}{C}{3 and 5}
\skill{distinguish observation from inference (M\&L p. 12).}
Statement 3 gives a \textbf{reason} (``because it gets more sun'') and statement 5 gives a \textbf{judgment} (``polluted''). Both go beyond what the student saw, so they are inferences. The others describe what was seen or measured.
\begin{whynot}
  \wrong{(A), (B), (D)}{statements 1, 2, and 4 are all observations (a color, a count, a measured temperature).}
\end{whynot}

\answer{2}{B}{2 and 4}
\skill{quantitative vs. qualitative data.}
Statement 2 is a \textbf{count} (14 and 3 frogs) and statement 4 is a \textbf{measurement} (18\,°C and 15\,°C). Quantitative data are ``numbers obtained by counting or measuring'' (M\&L p. 12).
\begin{whynot}
  \wrong{(A), (C)}{statement 1 is qualitative (a description); 3 and 5 are inferences, not data.}
  \wrong{(D)}{statement 3 is an inference.}
\end{whynot}

\answer{3}{B}{}
\skill{write a testable hypothesis.}
(B) names a cause (warmer water), a measurable result (number of frogs), and a comparison (warmed vs. unwarmed). It is written as an if \dots{} then prediction that an experiment can support or reject.
\begin{whynot}
  \wrong{(A)}{``happier'' can't be measured.}
  \wrong{(C)}{this is an inference stated as a conclusion, not a testable explanation. ``Polluted'' is also vague.}
  \wrong{(D)}{a value judgment (what \emph{should} happen) can't be tested by science.}
\end{whynot}

\answer{4}{B}{The color of the models}
\skill{identify variables.}
The researchers deliberately changed the model color, so it is the independent variable. The percentage attacked is what they measured: the \textbf{dependent} variable, choice (A).
\begin{whynot}
  \wrong{(C)}{the researchers didn't control which predators came.}
  \wrong{(D)}{time was held the same, so it was a controlled variable.}
\end{whynot}

\answer{5}{B}{2.6 times}
\skill{read a graph and calculate a ratio.}
Beach: dark models $\approx$ 73\% attacked, light models $\approx$ 28\%. 73 ÷ 28 $\approx$ \textbf{2.6}.
\begin{whynot}
  \wrong{(A)}{underestimates the ratio; (C) doesn't match the bar heights.}
  \wrong{(D)}{is the \emph{difference} in percentage points (73 -- 28 = 45), not ``how many times as often.''}
\end{whynot}

\answer{6}{B}{}
\skill{support a claim with data.}
On the beach, dark (non-matching) models were attacked more; inland, light (non-matching) models were attacked more. In both habitats, \textbf{mismatch = more attacks}.
\begin{whynot}
  \wrong{(A)}{false, because inland the dark models were attacked \emph{less} ($\approx$ 24\% vs. $\approx$ 78\%).}
  \wrong{(C)}{the experiment measures predation today. It can't prove how fur color \emph{evolved}, and science doesn't ``prove.''}
  \wrong{(D)}{overall, dark $\approx$ 73 + 24 = 97 and light $\approx$ 28 + 78 = 106, so dark models were not attacked more.}
\end{whynot}

\answer{7}{B}{}
\skill{evaluate experimental design.}
If only the beach were tested, a critic could say predators simply prefer dark objects. Testing a second habitat, where the \emph{opposite} color matches, shows that what matters is \textbf{matching the background}.
\begin{whynot}
  \wrong{(A)}{the number of predators wasn't the goal.}
  \wrong{(C)}{the independent variable (color) was kept.}
  \wrong{(D)}{experiments support hypotheses; they don't prove them.}
\end{whynot}

\answer{8}{B}{}
\skill{identify inductive reasoning.}
The class went from \textbf{many specific samples} to a \textbf{general rule}. That's induction. Its conclusion is only probable: sample 41 could be sterile (for example, a pond just treated with chemicals).
\begin{whynot}
  \wrong{(A), (C)}{deduction goes from general to specific, the other direction.}
  \wrong{(D)}{inductive conclusions are never guaranteed.}
\end{whynot}

\answer{9}{A}{}
\skill{identify deductive reasoning.}
(A) starts with a \textbf{general premise} (all insects have six legs) and applies it to a \textbf{specific case} (a beetle). If the premises are true, the conclusion must be true.
\begin{whynot}
  \wrong{(B), (C), (D)}{all generalize from past observations or a pattern, which is inductive (teacher's slides 4–5).}
\end{whynot}

\answer{10}{B}{}
\skill{describe a trend and relate it to a prediction.}
Mean distance fell from $\approx$ 2,400 km (0 advances) to $\approx$ 400 km (8 advances). Guesses got closer, as predicted.
\begin{whynot}
  \wrong{(A)}{the distance decreased, not increased.}
  \wrong{(C)}{data support predictions; they don't prove them, and a guess can still be wrong.}
  \wrong{(D)}{there is a large, steady change.}
\end{whynot}

\answer{11}{B}{}
\skill{sample size and reliability.}
``The larger the sample size, the more reliably researchers can analyze variation'' (M\&L p. 13). Combining class periods adds many more students per trial.
\begin{whynot}
  \wrong{(A)}{one student is a tiny sample.}
  \wrong{(C)}{searching the internet adds an uncontrolled variable and breaks the game's rules.}
  \wrong{(D)}{changing two variables at once makes it impossible to tell which one caused the result.}
\end{whynot}

\answer{12}{A}{}
\skill{use error bars to judge differences.}
From 0 to 5 to 10 g/m\textsuperscript{2} the bars rise and the error bars don't overlap, so those differences are likely real. The 10 and 20 g/m\textsuperscript{2} bars (58 and 60 cm) have \textbf{overlapping error bars}, so the data don't show a clear difference there. Growth may level off.
\begin{whynot}
  \wrong{(B)}{not every step is clearly different (see 10 vs. 20).}
  \wrong{(C)}{height clearly rose from 32 to 58 cm.}
  \wrong{(D)}{60 cm is not almost twice 58 cm.}
\end{whynot}

\answer{13}{B}{}
\skill{the role of a control group.}
The 0 g/m\textsuperscript{2} plots show \textbf{how grass grows without the treatment}, so any effect of nitrogen can be measured against them.
\begin{whynot}
  \wrong{(A)}{the dependent variable is grass height, not a group.}
  \wrong{(C)}{a control can't guarantee any result.}
  \wrong{(D)}{sunlight is kept the same by the design, not by the control group.}
\end{whynot}

\answer{14}{C}{}
\skill{the meaning of ``theory.''}
A scientific theory is ``a tested, highly-reliable scientific explanation \dots{} that unifies many repeated observations'' (M\&L p. 14). It can be revised when new evidence appears.
\begin{whynot}
  \wrong{(A)}{that is the everyday meaning, the opposite of the scientific one.}
  \wrong{(B)}{``theories do not become laws'' (M\&L p. 14). They are different kinds of statements.}
  \wrong{(D)}{scientists keep testing theories.}
\end{whynot}

\answer{15}{B}{}
\skill{the role of the scientific community.}
Reviewers ``look carefully for mistakes, oversights, unfair influences, or fraud'' (M\&L p. 18).
\begin{whynot}
  \wrong{(A)}{``Peer review does not guarantee that a piece of work is correct'' (M\&L p. 18).}
  \wrong{(C)}{reviewers are anonymous, independent experts, not friends.}
  \wrong{(D)}{replication is still needed.}
\end{whynot}

\answer{16}{C}{}
\skill{deliberate practice (the 1A target itself).}
Deliberate practice means using feedback to find the \textbf{exact gap}, practicing that skill, and checking again. A formative test is valuable \emph{because} it shows what to fix.
\begin{whynot}
  \wrong{(A)}{this wastes the feedback.}
  \wrong{(B)}{rereading everything is slow and doesn't target the gap.}
  \wrong{(D)}{by then it's too late to improve.}
\end{whynot}

\keyheading{Section II scoring guides and model answers}

\scoreheading{Question 1 (10 points)}

\begin{center}
\begin{tabularx}{\linewidth}{@{}llY@{}}
\toprule
Part & Points & Earn the point(s) for \dots{} \\
\midrule
(a) & 2 & Independent variable = hours of light per day (1). Dependent variable = seedling height after 10 days (1). \\
(b) & 2 & Two appropriate controlled variables, e.g., water, soil type, pot size, temperature, seed variety, light intensity or color (1). Explanation: if they differed, you couldn't tell whether light or the other variable caused the height differences (1). \\
(c) & 1 & The 4-hour group can serve as a baseline for comparing the other groups, \textbf{and} a limitation: without a 0-hour group you can't tell how seedlings grow with no light (for example, whether they grow tall and pale from stored food). \\
(d) & 2 & Correct axes and labels with units: x = hours of light per day, y = mean height (cm) (1). Four points (or bars) plotted correctly with error bars (1). \\
(e) & 1 & The trend: height increases as light increases from 4 to 12 hours, then levels off. \textbf{And} the 12 h and 16 h error bars overlap (5.2 $\pm$ 0.6 and 5.4 $\pm$ 0.7), so that difference is probably not real. \\
(f) & 1 & A claim supported by evidence \textbf{and} reasoning (see the model answer). \\
(g) & 1 & A prediction of about 5.2–5.8 cm, justified by the leveling-off trend (more light no longer adds much growth). \\
\bottomrule
\end{tabularx}
\end{center}

\modelanswer
\begin{parts}
  \item The independent variable is the number of hours of light per day. The dependent variable is the height of the seedlings after 10 days.
  \item The amount of water and the temperature should be the same for every group. If one group got more water or was kept warmer, it might grow taller for that reason, and the student couldn't tell whether light caused the difference.
  \item The 4-hour group can serve as a comparison group for the other three. Without a 0-hour group, the student can't say how much the seedlings would grow using only the food stored in the seed, so she can't measure the full effect of light.
  \item \emph{A line or bar graph: x-axis ``Hours of light per day'' (4, 8, 12, 16); y-axis ``Mean height (cm)'' from 0 to at least 6; points at 2.1, 3.8, 5.2, and 5.4, each with an error bar of the given size.}
  \item Height increases steeply from 4 to 12 hours of light (2.1 $\rightarrow$ 5.2 cm), then barely changes from 12 to 16 hours (5.2 $\rightarrow$ 5.4 cm). The error bars for 12 and 16 hours overlap (4.6–5.8 and 4.7–6.1 cm), so the difference between those two groups is probably not real.
  \item \textbf{Claim:} more light increases radish seedling growth, but only up to about 12 hours per day. \textbf{Evidence:} mean height rose from 2.1 cm at 4 hours to 5.2 cm at 12 hours, with no overlap in those error bars, but 16 hours (5.4 cm) was not clearly different from 12 hours. \textbf{Reasoning:} plants use light energy for photosynthesis to build the molecules they need to grow. Beyond about 12 hours, something else (such as water, nutrients, or the plant's growth rate) likely limits growth.
  \item About 5.4 cm (anything from 5.2 to 5.8 cm). The data level off after 12 hours, so 20 hours probably won't add much more growth.
\end{parts}

\scoreheading{Question 2 (4 points)}

\begin{center}
\begin{tabularx}{\linewidth}{@{}lY@{}}
\toprule
Part & Earn 1 point for \dots{} \\
\midrule
(a) & One observation (dead fish downstream; no dead fish upstream; the factory is between them) \textbf{and} one inference (e.g., something from the factory killed the fish). \\
(b) & Inductive, because it reasons from specific observations to a general cause. It is not certain because another cause could explain the same observations. \\
(c) & A testable if \dots{} then hypothesis \textbf{and} the result that would support it. \\
(d) & An observation that points to a different cause or contradicts the conclusion. \\
\bottomrule
\end{tabularx}
\end{center}

\modelanswer
\begin{parts}
  \item Observation: there are dead fish downstream from the factory but none upstream. Inference: something released by the factory killed the fish.
  \item This is inductive reasoning: the biologist went from specific observations to a general explanation. It isn't certain, because another cause could produce the same pattern, such as a disease, low oxygen, or a different source of pollution downstream.
  \item \emph{If} the factory's wastewater is toxic to fish, \emph{then} fish kept in tanks of river water collected just below the factory's pipe will die faster than fish kept in water collected upstream (with temperature, oxygen, and fish size kept the same). Supporting result: the fish in the downstream water die sooner.
  \item Finding dead fish in a stream that the factory's water never reaches, or finding that the water below the pipe has normal oxygen and no toxic chemicals.
\end{parts}

\scoreheading{Question 3 (4 points)}

\begin{center}
\begin{tabularx}{\linewidth}{@{}lY@{}}
\toprule
Part & Earn 1 point for \dots{} \\
\midrule
(a) & The dark models (camouflaged on dark soil). \\
(b) & About 54 percentage points ($\approx$ 78\% -- $\approx$ 24\%). Accept 50–58. \\
(c) & A valid reason: models are identical except for color, and they don't differ in behavior, smell, health, or movement, so color is the only difference. \\
(d) & The camouflaged model in each habitat was attacked less, which is what the hypothesis predicts. \\
\bottomrule
\end{tabularx}
\end{center}

\modelanswer
\begin{parts}
  \item The dark models are the control group in the inland habitat, because they match the dark soil (camouflaged).
  \item Light models $\approx$ 78\% attacked, dark models $\approx$ 24\%: a difference of about \textbf{54 percentage points}.
  \item Real mice differ in many ways besides color: how they move, hide, smell, and behave. Identical models remove all of those differences, so color is the only variable, and any difference in attacks can be linked to color.
  \item The hypothesis predicts that mice matching their surroundings are attacked less. In both habitats, the matching models were attacked far less ($\approx$ 28\% vs. $\approx$ 73\% on the beach; $\approx$ 24\% vs. $\approx$ 78\% inland). Because the result held when the matching color was reversed, it supports camouflage rather than a simple preference for one color.
\end{parts}

\scoreheading{Question 4 (4 points)}

\begin{center}
\begin{tabularx}{\linewidth}{@{}lY@{}}
\toprule
Part & Earn 1 point for \dots{} \\
\midrule
(a) & Independent experts in the same field review the paper's methods, data, and reasoning for mistakes, bias, or fraud before it is published. \\
(b) & Science doesn't prove ideas; evidence supports them. Reviewers can miss errors, and new evidence can change conclusions. \\
(c) & Other scientists repeat the study; getting the same result makes it unlikely the first result was chance, error, or bias. \\
(d) & Both use outside feedback to find and fix weaknesses before a final judgment. \\
\bottomrule
\end{tabularx}
\end{center}

\modelanswer
\begin{parts}
  \item Before publication, anonymous, independent scientists who work in the same field read the paper. They check the methods, data, and reasoning for mistakes, oversights, bias, or fraud, and the authors must fix problems before the paper is accepted.
  \item Peer review only shows that the work meets the community's standards. Reviewers can miss errors, and science never proves an idea beyond all doubt: a hypothesis is supported by evidence, and future evidence could change the conclusion.
  \item If other scientists repeat the experiment and get the same result, the result is unlikely to be due to chance, a mistake, or one group's bias. Each successful replication increases confidence in the conclusion.
  \item Formative assessments in class are like peer review: someone else checks your work and points out mistakes \emph{before} the final (summative) judgment, so you can fix the gaps. In both cases, feedback improves the final product.
\end{parts}

\end{document}
